% Options for packages loaded elsewhere
\PassOptionsToPackage{unicode}{hyperref}
\PassOptionsToPackage{hyphens}{url}
\documentclass[
]{article}
\usepackage{xcolor}
\usepackage{amsmath,amssymb}
\setcounter{secnumdepth}{-\maxdimen} % remove section numbering
\usepackage{iftex}
\ifPDFTeX
  \usepackage[T1]{fontenc}
  \usepackage[utf8]{inputenc}
  \usepackage{textcomp} % provide euro and other symbols
\else % if luatex or xetex
  \usepackage{unicode-math} % this also loads fontspec
  \defaultfontfeatures{Scale=MatchLowercase}
  \defaultfontfeatures[\rmfamily]{Ligatures=TeX,Scale=1}
\fi
\usepackage{lmodern}
\ifPDFTeX\else
  % xetex/luatex font selection
\fi
% Use upquote if available, for straight quotes in verbatim environments
\IfFileExists{upquote.sty}{\usepackage{upquote}}{}
\IfFileExists{microtype.sty}{% use microtype if available
  \usepackage[]{microtype}
  \UseMicrotypeSet[protrusion]{basicmath} % disable protrusion for tt fonts
}{}
\makeatletter
\@ifundefined{KOMAClassName}{% if non-KOMA class
  \IfFileExists{parskip.sty}{%
    \usepackage{parskip}
  }{% else
    \setlength{\parindent}{0pt}
    \setlength{\parskip}{6pt plus 2pt minus 1pt}}
}{% if KOMA class
  \KOMAoptions{parskip=half}}
\makeatother
\setlength{\emergencystretch}{3em} % prevent overfull lines
\providecommand{\tightlist}{%
  \setlength{\itemsep}{0pt}\setlength{\parskip}{0pt}}
\usepackage{bookmark}
\IfFileExists{xurl.sty}{\usepackage{xurl}}{} % add URL line breaks if available
\urlstyle{same}
% fallback for those not using the hyperref driver hyperxmp:
\makeatletter
\@ifundefined{xmpquote}{\newcommand{\xmpquote}[1]{#1}}{}
\makeatother
\hypersetup{
  hidelinks,
  pdfcreator={LaTeX via pandoc}}

\author{}
\date{}

\begin{document}

{
\setcounter{tocdepth}{3}
\tableofcontents
}
\section{App Store Distribution Exception
Plan}\label{app-store-distribution-exception-plan}

\subsection{Purpose}\label{purpose}

Apple platform distribution terms may create a distribution problem for
a GPL-licensed application. This document plans the review needed to
decide whether a narrowly scoped distribution exception is appropriate
for Password Bonobo. It does not decide that question and does not
publish exception language.

\subsection{Proposed scope
constraints}\label{proposed-scope-constraints}

Any future exception would be limited to Bonobo-authored code for which
every relevant contributor granted the necessary permission. It would
not apply to Gorilla-derived code, which is excluded from Bonobo product
builds, or to dependencies whose licenses or terms are incompatible with
the approved iOS distribution model.

No contributor permission, license exception, or iOS distribution right
is created by this planning document.

\subsection{Required reviews}\label{required-reviews}

Before any decision, the project must complete and record these reviews:

\begin{itemize}
\tightlist
\item
  A dependency review covering licenses, distribution terms, and iOS
  build inclusion.
\item
  A contributor-rights review covering the permission needed for the
  final approved scope.
\item
  An Apple-term review using the terms applicable at the time of the
  decision.
\item
  A legal review of the distribution model, proposed scope, and final
  wording.
\end{itemize}

\subsection{Decision gates}\label{decision-gates}

External contributions remain closed until the contribution terms can
preserve GPL-3.0-or-later licensing and any approved distribution
exception rights for Bonobo-authored iOS code. The repository owner must
deliberately publish the contribution terms before accepting external
contributions.

Before the first iOS distribution build, the project must confirm the
final exception decision, contributor permissions, dependency
compatibility, Apple-term review, and legal review. The iOS build must
exclude Gorilla-derived code and any incompatible dependency.

\subsection{Status and limitations}\label{status-and-limitations}

This is a planning document, not license text or legal advice. It must
not be cited as final exception language, contributor permission, or
approval to distribute an iOS build.

\end{document}
